
\documentclass[12pt,a4paper,UTF8]{ctexart}
\usepackage{geometry}
\geometry{a4paper,left=2.5cm,right=2.5cm,top=3cm,bottom=3cm}
\usepackage{graphicx}
\usepackage{hyperref}
\usepackage{enumitem}
\usepackage{tabularx}
\usepackage{booktabs}
\usepackage{longtable}
\usepackage{caption}
\usepackage{fancyhdr}
\usepackage{xcolor}
\usepackage{footmisc}
\usepackage{fontspec}
\usepackage{listings}

% 使用TeX内置的Fandol中文字体
\setCJKmainfont[BoldFont=FandolSong-Bold,ItalicFont=FandolKai]{FandolSong}
\setCJKsansfont{FandolHei}
\setCJKmonofont{FandolFang}
\setmainfont{FandolSong}

% 页眉页脚
\pagestyle{fancy}
\fancyhf{}
\fancyhead[L]{\leftmark}
\fancyhead[R]{\thepage}
\fancyfoot[C]{黔西南自驾攻略 v1.0}

% 链接颜色
\hypersetup{
    colorlinks=true,
    linkcolor=blue,
    filecolor=magenta,
    urlcolor=cyan,
}

% 标题设置
\title{\Huge \textbf{🏔️ 黔西南7天自驾全攻略}}
\author{家庭出游指南}
\date{2026年5月}

\begin{document}

\maketitle
\thispagestyle{empty}

\newpage
\tableofcontents
\newpage


\section{行程总览}

\begin{verbatim}
Day 0：各地 → 贵阳（集合取车）
Day 1：贵阳 → 织金洞 → 盘州
Day 2：盘州 → 乌蒙大草原（滑翔伞）→ 北盘江大桥 → 兴义
Day 3：兴义 → 万峰林（热气球）
Day 4：兴义 → 马岭河峡谷（漂流）→ 万峰湖
Day 5：兴义 → 贞丰双乳峰 → 格凸河（攀岩）
Day 6：格凸河 → 坝陵河大桥（蹦极/跳伞）→ 平塘天眼
Day 7：平塘 → 贵阳 → 各地
\end{verbatim}

\subsection{基本信息}
\begin{itemize}
\item 总里程：约1500km
\item 日均车程：3-4小时
\item 人均预算：约5000元（基础版）
\item 最佳时间：7-8月（避暑天堂，18-28度）
\item 适合人数：5-8人（2台车）
\item 难度：轻松休闲
\end{itemize}

\subsection{方案亮点}
\begin{itemize}
\item ✅ 路线完美：全程环线，几乎没有回头路，每天车程3-4小时不累
\item ✅ 体验分层：每个刺激项目都有分流方案，没有人需要干等
\item ✅ 景点顶级：6个国家地理级别的景观
\item ✅ 游客极少：大部分时候都是包场体验
\item ✅ 避暑天堂：7-8月全国最热的时候，贵州只有18-28度
\end{itemize}

\newpage

\section{每日详细行程}

\subsection{Day 0：各地 → 贵阳}
\begin{itemize}
\item 各地航班陆续到达贵阳龙洞堡机场
\item 机场租车点集合，取车，验车
\item 晚餐：老凯俚酸汤鱼
\item 晚上车队准备会，分发手台
\end{itemize}

\subsubsection{租车建议}
\begin{itemize}
\item 平台：一嗨/神州机场店
\item 车型：5-6人选1台7座MPV+1台轿车，7-8人选2台SUV
\item 保险：一定要买全额保险，三者险至少200万
\item 取车：拍好所有划痕、轮胎、玻璃的照片视频
\end{itemize}

\subsection{Day 1：贵阳 → 织金洞 → 盘州}
\begin{itemize}
\item 08:00 酒店早餐，退房出发
\item 10:00-13:30 织金洞游览（3.5小时）
\item 织金洞："黄山归来不看岳，织金洞外无洞天"
\item 洞内16度，带薄外套，不要用手摸钟乳石
\item 14:30 出发去盘州
\item 17:00 到达盘州，入住酒店
\end{itemize}

\subsection{Day 2：盘州 → 乌蒙大草原 → 北盘江大桥 → 兴义}
\begin{itemize}
\item 07:30 出发，赶云海
\item 08:30-12:00 乌蒙大草原游玩
\item 海拔2800米，温度比山下低10度，带厚外套
\item 不要开车压草坪，会被罚款500
\item 15:30 北盘江大桥拍照（世界第一高桥，565米）
\item 18:30 到达兴义
\end{itemize}

\subsection{Day 3：万峰林全天}
\begin{itemize}
\item 上午：景区观光车，山上看全景
\item 下午：租电动车在峰林骑行（30元/天）
\item 推荐路线：将军峰→下纳灰→中纳灰→上纳灰
\item 不要进任何收费的"村寨景点"，都是坑
\end{itemize}

\subsection{Day 4：马岭河峡谷 + 万峰湖}
\begin{itemize}
\item 上午：马岭河峡谷漂流或徒步
\item 漂流：239元/人，22公里，2.5小时
\item 徒步：3公里平路栈道，看十几条瀑布
\item 下午：万峰湖坐船
\end{itemize}

\subsection{Day 5：兴义 → 贞丰 → 格凸河}
\begin{itemize}
\item 双乳峰：路边打卡即可，不用进景区
\item 下午到达格凸河
\item 16:00 蜘蛛人表演（每天只有一场，别错过）
\item 下午2-4点能看到神光
\end{itemize}

\subsection{Day 6：格凸河 → 坝陵河大桥 → 平塘天眼}
\begin{itemize}
\item 坝陵河大桥：国内最高蹦极370米，2000元/人
\item 不玩蹦极的可以走观光栈道，100元/人
\item 平塘天眼：所有电子设备必须寄存，可以带胶卷相机
\item 观景台爬789级台阶
\end{itemize}

\subsection{Day 7：平塘 → 贵阳 → 各地}
\begin{itemize}
\item 上午出发回贵阳
\item 午餐后采购特产
\item 机场还车，返程
\end{itemize}

\newpage

\section{刺激项目分流方案}

所有刺激项目都设计了分流体验，两边同时进行，互不等待。

\subsection{马岭河峡谷漂流}
\begin{itemize}
\item 价格：239元/人
\item 玩的人：漂流起点出发，2.5小时漂完全程
\item 不玩的人：走峡谷徒步栈道，看瀑布，同时在终点汇合
\item 注意：穿凉鞋，不要穿拖鞋（会冲跑）
\end{itemize}

\subsection{坝陵河大桥蹦极}
\begin{itemize}
\item 价格：2000元/人
\item 玩的人：蹦极/高空秋千/跳伞
\item 不玩的人：走观光栈道，看别人蹦极，围观也很刺激
\item 提前一天预约
\end{itemize}

\subsection{乌蒙大草原滑翔伞}
\begin{itemize}
\item 价格：680元/人
\item 飞行时间：10-15分钟
\item 不玩的人：逛天池，骑马，看滑翔伞起飞降落
\end{itemize}

\subsection{格凸河攀岩体验}
\begin{itemize}
\item 价格：250元/人
\item 初级体验线，普通人也能爬
\item 不玩的人：看蜘蛛人表演，坐船游河，穿洞看神光
\end{itemize}

\subsection{万峰林热气球}
\begin{itemize}
\item 价格：180元/人
\item 系留飞，30分钟
\item 不玩的人：骑电动车逛峰林
\end{itemize}

\newpage

\section{住宿推荐}

\begin{table}[h]
\centering
\caption{住宿总览}
\begin{tabular}{llcl}
\toprule
日期 & 酒店 & 价格/间 & 说明 \\
\midrule
Day 0 & 维也纳酒店（机场店） & 200 & 离机场近 \\
Day 1 & 盘州国际花园 & 250 & 盘州最好 \\
Day 2-4 & 兴义富康国际 & 280 & \textbf{连住3晚！} \\
Day 5 & 格凸河镇酒店 & 200 & 比景区便宜 \\
Day 6 & 平塘天文小镇 & 250 & 离天眼近 \\
\bottomrule
\end{tabular}
\end{table}

\subsection{订房Tips}
\begin{itemize}
\item 暑期提前2周订，不然涨价甚至没房
\item 三个平台比价：携程、美团、飞猪
\item 选可以免费取消的
\item 所有推荐酒店都有免费停车场
\end{itemize}

\newpage

\section{预算明细}

\subsection{基础版人均预算}
\begin{itemize}
\item 往返大交通：1500元
\item 租车+油费+过路费：600元
\item 住宿（6晚）：700元
\item 餐饮（7天）：700元
\item 门票：800元
\item 其他杂项：400元
\item \textbf{总计：约5000元/人}
\end{itemize}

\subsection{省钱技巧}
\begin{itemize}
\item 周二周三订票最便宜
\item 提前2周订酒店，便宜30%
\item 美团买门票比现场便宜
\item 不要在景区门口吃饭
\end{itemize}

\newpage

\section{行前准备清单}

\subsection{车队装备}
\begin{itemize}
\item 对讲机2台（山区没信号）
\item 车载充电器、手机支架
\item 拖车绳、搭电线
\end{itemize}

\subsection{衣物}
\begin{itemize}
\item ✅ 速干衣裤（漂流玩水用）
\item ✅ 凉鞋/溯溪鞋（必带！不要穿拖鞋）
\item ✅ 薄外套（洞里16度）
\item ✅ 厚外套/冲锋衣（乌蒙15度）
\item ✅ 雨衣（比雨伞好用）
\item ✅ 帽子、墨镜
\end{itemize}

\subsection{药品}
\begin{itemize}
\item ✅ 肠胃药（吃辣必备）
\item ✅ 感冒药
\item ✅ 创可贴、碘伏
\item ✅ 晕车药（盘山路）
\item ✅ 驱蚊液
\end{itemize}

\subsection{电子设备}
\begin{itemize}
\item ✅ 手机、充电器
\item ✅ 充电宝（20000mAh以内）
\item ✅ 防水袋（漂流用）
\end{itemize}

\subsection{特别提醒}
\begin{itemize}
\item 贵州7-8月是雨季，但都是阵雨，下下停停
\item 贵州夏天真的不热，25度左右，晚上要盖被子
\item 紫外线很强，防晒霜、帽子、墨镜缺一不可
\end{itemize}

\newpage

\section{车队自驾规则}

\subsection{基本规则}
\begin{enumerate}
\item 头车尾车制度：1号车头车带路，2号车尾车收尾
\item 每15分钟手台报一次位置
\item 高速保持100-200米车距，贵州限速110
\item 每2小时进服务区休息一次，换人开车
\item 统一行动：加油、休息、吃饭一起
\end{enumerate}

\subsection{贵州路况特点}
\begin{itemize}
\item 高速多隧道多桥梁，进隧道提前减速开灯
\item 限速严，区间测速多
\item 7-8月多雨多雾，能见度低就近休息
\item 盘山路弯多路窄，弯道鸣笛靠右行，不要超车
\end{itemize}

\subsection{应急预案}
\begin{itemize}
\item 两车走散了：约定下一个服务区集合，不要在高速停车
\item 车辆故障：所有人撤到护栏外，打12122高速救援
\item 事故：先看人有没有事，再拍照，报警，报保险
\item 暴雨大雾：减速到60以下，就近进服务区
\end{itemize}

\subsection{重要电话}
\begin{itemize}
\item 高速救援：12122
\item 交通事故：122
\item 急救：120
\end{itemize}

\newpage

\section{景点避坑指南}

\subsection{织金洞}
❌ 不要跟团，旅行团只走一半就催你出来
❌ 不要请导游，跟着别的团听免费讲解
✅ 带薄外套，洞里16度

\subsection{乌蒙大草原}
❌ 不要开车压草坪，罚款500不讲价
❌ 不要在景区门口吃烤全羊
✅ 早上去，大概率能看到云海

\subsection{万峰林}
❌ 不要进任何收费"村寨景点"，都是坑
❌ 不要在景区门口吃蛋炒饭，去将军峰
✅ 租电动车30元骑一下午

\subsection{马岭河峡谷}
❌ 不要穿拖鞋漂流，一定会冲跑
❌ 不要坐观光电梯，自己走10分钟，纯坑

\subsection{格凸河}
❌ 不要上午去，上午看不到神光
❌ 不要错过下午4点的蜘蛛人表演，每天只有一场
✅ 住镇上，不要住景区里

\subsection{坝陵河大桥}
❌ 不要导航"坝陵河景区"，导航"坝陵河大桥研学基地"
❌ 不要在高速上停车拍照，扣6分

\subsection{平塘天眼}
❌ 不要抱侥幸带手机，过安检门，查到必须寄存
✅ 可以带胶卷相机，景区有卖200元一个

\newpage

\section{实用信息}

\subsection{必备APP}
\begin{itemize}
\item 导航：高德地图（贵州数据比百度全）
\item 住宿：携程、美团、飞猪（三个比价）
\item 门票：美团（比现场便宜）
\item 路况：贵州高速交警微博
\end{itemize}

\subsection{贵州美食推荐}
\subsubsection{贵阳}
\begin{itemize}
\item 酸汤鱼：老凯俚酸汤鱼
\item 丝娃娃、肠旺面、烙锅
\end{itemize}

\subsubsection{兴义}
\begin{itemize}
\item 蛋炒饭：万峰林将军峰
\item 刷把头、鸡肉汤圆、羊肉粉
\end{itemize}

\subsection{贵州特产}
\begin{itemize}
\item 茅台酒（机场免税店可能有）
\item 都匀毛尖（茶叶）
\item 刺梨干（酸甜，富含VC）
\item 老干妈
\end{itemize}

\subsection{最后提醒三遍}
\begin{center}
\Large \textbf{带厚外套！带厚外套！带厚外套！}
\end{center}
\begin{center}
\Large \textbf{带凉鞋！带凉鞋！带凉鞋！}
\end{center}
\begin{center}
\Large \textbf{带雨衣！带雨衣！带雨衣！}
\end{center}

\vspace{2cm}
\begin{center}
\Huge 祝大家玩得开心！🚗🏔️🌊
\end{center}

\end{document}
